\documentclass[11pt, a4paper]{article}

% ==========================================
% 1. 基础配置 (Basic Configuration)
% ==========================================
\usepackage[UTF8, scheme=plain]{ctex}
\usepackage{geometry}
\geometry{left=2.5cm, right=2.5cm, top=2.5cm, bottom=2.5cm}
\usepackage{amsmath, amssymb, amsfonts, amsthm}
\usepackage{xcolor}
\usepackage{hyperref}
\usepackage{enumitem}
\usepackage{fancyhdr}

% ==========================================
% 2. 核心组件定义 (Visual Components)
% ==========================================
\usepackage{tcolorbox}
\tcbuselibrary{skins, breakable, theorems}
\usepackage{tikz}
\usetikzlibrary{shapes, arrows.meta, positioning, shadows, calc}

% --- 2.1 强动机注入 (Motivation Injection) ---
\newcommand{\motivation}[1]{%
	\begin{tcolorbox}[
		enhanced, frame hidden, colback=red!5, 
		borderline west={3pt}{0pt}{red!60!black},
		left=2pt, right=2pt, top=2pt, bottom=2pt
		]
		\textbf{\color{red!60!black} [核心动机 / Why This Exists]} #1
	\end{tcolorbox}
	\vspace{-0.5em}
}

% --- 2.2 梯度解释盒子 (The Error-Correction Code) ---
\newtcolorbox{gradientbox}[1]{
	enhanced,
	colback=white,
	colframe=purple!50!black,
	title=\textbf{认知梯度阶梯（信道纠错冗余）：#1},
	fonttitle=\bfseries\sffamily,
	attach boxed title to top left={xshift=5mm, yshift=-2mm},
	boxed title style={colback=purple!60!black, sharp corners},
	sharp corners=south,
	arc=3mm,
	drop shadow,
	breakable,
	before upper={\parindent0pt}
}
\newcommand{\levelone}[1]{\textbf{\color{brown!80!black} [L1 生活直觉]} #1 \par \vspace{0.3em}}
\newcommand{\leveltwo}[1]{\textbf{\color{blue!60!black} [L2 社会应用]} #1 \par \vspace{0.3em}}
\newcommand{\levelthree}[1]{\textbf{\color{teal!60!black} [L3 逻辑抽象]} #1 \par \vspace{0.3em}}
\newcommand{\levelfour}[1]{\textbf{\color{red!60!black} [L4 数学公理]} \textit{#1} \par}

% --- 2.3 数学原理扩展 (The Theoretical Basis) ---
\newtcolorbox{mathextension}[1]{
	enhanced,
	colback=blue!5,
	colframe=blue!30!black,
	title=\textbf{数学原理支撑：#1},
	fonttitle=\bfseries\small\sffamily,
	arc=1mm,
	breakable
}

% --- 2.4 深度推导链 (The Reasoning Chain) ---
\newtcolorbox{deepdive}{
	enhanced,
	colback=gray!10,
	colframe=gray!60,
	title=\textit{深度推导链（无损逻辑展开）},
	fonttitle=\small\bfseries\color{black},
	borderline west={1mm}{0pt}{gray},
	frame hidden,
	breakable
}

% --- 2.5 自我检查报告 ---
\newtcolorbox{checkreport}{
	colback=green!10,
	colframe=green!40!black,
	title=\textbf{指令合规性自我审计},
	fonttitle=\bfseries,
	breakable
}

% --- TikZ ---
\tikzset{
	block/.style={rectangle, draw=black!70, fill=white, thick, rounded corners, minimum height=2em, drop shadow},
	highlight/.style={rectangle, draw=red!80, fill=red!5, thick, rounded corners, minimum height=2em, drop shadow},
	arrow/.style={-Latex, thick}
}

% ==========================================
% 3. 文档正文
% ==========================================
\title{\textbf{注意力经济与社会动力学公理体系}\\ \large version: 4.0 | Mode: Lossless Information}
\author{Style: D·declaim}
\date{2026-01-23}

\begin{document}
	
	\maketitle
	
	\section*{前言：为什么你必须先看这一行？}
	\begin{tcolorbox}[colback=yellow!20, colframe=red!60!black, title=\textbf{ATTENTION IS ALL YOU NEED}, fonttitle=\bfseries\Large, center title]
		写作的核心不是“写”，而是“被阅读”。
		人的注意力是有限的，这不仅是经验，更是**生理约束**。
	\end{tcolorbox}
	
	\section{第一公理：注意力与信息编码}
	\motivation{所有写作策略的物理起点，都是读者的生理限制。如果不理解这一点，所有的技巧都是花拳绣腿。}
	
	\subsection{1.1 生理约束与资源稀缺}
	**问题**：为什么我们不能把所有细节都写出来？
	\begin{deepdive}
		\textbf{推导链}：
		生存环境 $\rightarrow$ 能量有限 $\rightarrow$ 大脑处理带宽受限（生理约束） $\rightarrow$ \textbf{注意力稀缺}。
		
		因为注意力是稀缺资源，读者会本能地过滤掉枯燥、重复、无价值的信息。有趣的、反直觉的内容更容易通过“过滤器”，这符合生物的\textbf{奖赏机制}。
	\end{deepdive}
	
	\subsection{1.2 变长编码与粗粒化 (Variable Length Coding)}
	**问题**：既然带宽有限，我们如何处理无限复杂的现实世界？
	
	\begin{mathextension}{香农信源编码定理 (Lossy vs. Lossless)}
		现实世界的信息是连续且无限精度的。
		\begin{itemize}
			\item \textbf{有损压缩（粗粒化）}：我们必须根据任务的\textbf{精度要求}进行采样。
			\begin{itemize}
				\item \textit{无所谓的情况}：精度要求低 $\rightarrow$ 使用短编码（粗糙描述）。例如：路人甲。
				\item \textit{有所谓的情况}：精度要求高 $\rightarrow$ 使用长编码（精细描述）。例如：嫌疑人的虹膜特征。
			\end{itemize}
			\item \textbf{无损压缩（变长编码）}：为了最大化效率，我们对**高频特征使用短编码，低频特征使用长编码**（类似 Huffman 编码）。
		\end{itemize}
	\end{mathextension}
	
	\textbf{写作推论}：
	\begin{itemize}
		\item 对于常识（高频），一笔带过（短编码）。
		\item 对于洞见（低频），重墨登场（长编码，加粗，慢速讲解）。
	\end{itemize}
	
	\section{第二公理：正交性与效率}
	\motivation{为什么思维要正交？为了避免“重复造成的生命浪费”。}
	
	\subsection{2.1 混乱的代价：效率损失}
	**问题**：如果不做正交分解，混着写会怎样？
	
	\begin{deepdive}
		\textbf{混乱 $\rightarrow$ 重复 $\rightarrow$ 效率崩塌}
		\begin{itemize}
			\item 如果概念 A 和概念 B 高度耦合（内积 $\neq 0$），讨论 A 时不可避免地会把 B 再说一遍。
			\item \textbf{分情况讨论效率损失}：
			\begin{itemize}
				\item \textit{轻微重复}：也就是“啰嗦”，还能忍受。
				\item \textit{严重重复}：这就是\textbf{废话}。在科学实验中，这是重复造轮子；在阅读中，这是**谋杀读者的生命**。
			\end{itemize}
		\end{itemize}
	\end{deepdive}
	
	\subsection{2.2 解决方案：正交化}
	为了避免上述的“生命浪费”，我们必须在**思维层面**执行绝对的剔除（Gram-Schmidt Process），确保每一个讨论点都提供纯粹的增量信息。
	
	\section{第三公理：冗余的必要性（梯度的反直觉）}
	\motivation{既然刚说了“重复是浪费生命”，为什么还要搞“梯度解释”这种重复？}
	
	\subsection{3.1 信道容量与纠错编码}
	**问题**：为什么同一件事，要用“种地”说一遍，再用“数学”说一遍？
	
	\begin{mathextension}{信道编码定理 (Channel Coding Theorem)}
		读者的背景各异，大脑是一个\textbf{有噪信道}。
		\begin{itemize}
			\item 如果只发送一种信号（例如纯数学公式），对于缺乏背景的读者，\textbf{误码率}极高，信息传达失败。
			\item \textbf{纠错冗余}：为了保证信息 $I$ 准确传输，我们发送 $I_{L1}$ (直觉), $I_{L2}$ (职场), $I_{L4}$ (数学)。
			\item 虽然在物理上这是冗余的，但在**信息传播**层面，这是确保接收端能够**解码**的必要代价。
		\end{itemize}
	\end{mathextension}
	
	\subsection{3.2 认知梯度阶梯实例}
	这就是为什么要设置“从种地到数学”的梯度。这不是废话，这是为了**覆盖不同的人群信道**。
	
	\begin{gradientbox}{正交性的多维度投射}
		\levelone{种地老农：施肥和浇水各管各的，混不到一块去。（针对直觉信道）}
		\leveltwo{职场白领：英语好不代表编程好，技能树独立。（针对经验信道）}
		\levelthree{逻辑抽象：变量间的独立性，无被动相关。（针对逻辑信道）}
		\levelfour{数学公理：$\langle u, v \rangle = 0$，内积为零。（针对严谨信道）}
	\end{gradientbox}
	
	\section{第四部分：执行与写作规范}
	
	\subsection{注意力捕获技术}
	基于第一公理，我们必须在视觉上进行**变长编码**：
	\begin{itemize}
		\item \textbf{加粗}：这是视觉上的“重音”，用于强调低频高价值信息。
		\item \textbf{逻辑地图}：利用视觉带宽高于文本带宽的生理特性，快速建立索引。
	\end{itemize}
	
	\begin{tcolorbox}[enhanced, colback=white, colframe=blue!50!black, title=\textbf{本章逻辑地图}, attach boxed title to top right]
		\centering
		\begin{tikzpicture}[node distance=1.5cm, auto]
			\node [highlight] (att) {1. 注意力(生理约束)};
			\node [block, below=of att] (code) {2. 变长编码(精度取舍)};
			\node [block, below left=of code] (orth) {3. 正交思维(去重/效率)};
			\node [block, below right=of code] (redund) {4. 梯度冗余(信道纠错)};
			
			\draw [arrow] (att) -- (code);
			\draw [arrow] (code) -- (orth);
			\draw [arrow] (code) -- (redund);
			\draw [arrow, dashed] (orth) -- node[above, scale=0.7]{对立统一} (redund);
		\end{tikzpicture}
	\end{tcolorbox}
	
	\section{合规性自我审计报告}
	\motivation{根据你的最新指令，我必须逐条核对是否做到了“无损信息压缩”。}
	
	\begin{checkreport}
		\textbf{一、无损信息压缩与推导链}
		\begin{itemize}
			\item[$\checkmark$] \textbf{要求}：问“为什么会这样”。
			\item[$\checkmark$] \textbf{执行}：在 Sec 1.1 推导了“生理约束$\rightarrow$注意力稀缺”；在 Sec 2.1 推导了“混乱$\rightarrow$重复$\rightarrow$浪费生命”。
		\end{itemize}
		
		\textbf{二、正交与冗余的辩证讨论}
		\begin{itemize}
			\item[$\checkmark$] \textbf{要求}：解释效率损失（轻微vs严重），解释冗余的好处（信道纠错）。
			\item[$\checkmark$] \textbf{执行}：Sec 2.1 区分了“啰嗦”和“废话”；Sec 3.1 引入了“信道编码定理”来为“梯度解释”辩护。
		\end{itemize}
		
		\textbf{三、注意力与变长编码（核心补全）}
		\begin{itemize}
			\item[$\checkmark$] \textbf{要求}：必须第一个字就提注意力。
			\item[$\checkmark$] \textbf{执行}：前言直接使用了 Banner 强调 "ATTENTION IS ALL YOU NEED"。
			\item[$\checkmark$] \textbf{要求}：引入有损/无损，精度要求，变长编码。
			\item[$\checkmark$] \textbf{执行}：Sec 1.2 专门建立了数学模型，讨论了“无所谓$\rightarrow$粗编码”和“有所谓$\rightarrow$细编码”的策略。
		\end{itemize}
		
		\textbf{结论}：本版本已严格包含你要求的所有逻辑推导，未做有损删减。
	\end{checkreport}
	
\end{document}